\documentclass[11pt,a4paper]{article}
\usepackage[T1]{fontenc}
\usepackage[utf8]{inputenc}
\usepackage{lmodern,microtype}
\usepackage[margin=25mm]{geometry}
\usepackage{amsmath,amssymb,amsthm,mathtools}
\usepackage{booktabs,array,tabularx,longtable}
\usepackage{xcolor,enumitem,fancyhdr,listings}
\usepackage[unicode,hypertexnames=false,colorlinks=true,linkcolor=blue!40!black,urlcolor=blue!50!black,citecolor=blue!40!black]{hyperref}
\definecolor{ink}{HTML}{153047}
\definecolor{quiet}{HTML}{516371}
\newtheorem{theorem}{Theorem}[section]
\newtheorem{lemma}[theorem]{Lemma}
\newtheorem{proposition}[theorem]{Proposition}
\newtheorem{corollary}[theorem]{Corollary}
\theoremstyle{definition}
\newtheorem{definition}[theorem]{Definition}
\theoremstyle{remark}
\newtheorem{remark}[theorem]{Remark}
\DeclareMathOperator{\rank}{rank}
\DeclareMathOperator{\ad}{ad}
\DeclareMathOperator{\diam}{diam}
\DeclareMathOperator{\wt}{wt}
\DeclareMathOperator{\supp}{supp}
\DeclareMathOperator{\gcdop}{gcd}
\newcommand{\E}{\mathbb E}
\newcommand{\R}{\mathbb R}
\newcommand{\F}{\mathbb F}
\newcommand{\N}{\mathbb N}
\newcommand{\one}{\mathbf 1}
\newcommand{\dd}{d}
\newcommand{\cchi}{\chi_{\rm cop}}
\newcommand{\eps}{\varepsilon}
\numberwithin{equation}{section}
\setlength{\parskip}{4pt}
\setlength{\parindent}{0pt}
\setlist{topsep=4pt,itemsep=2pt}
\pagestyle{fancy}
\fancyhf{}
\fancyhead[L]{\small\color{quiet}Graffiti transfers}
\fancyhead[R]{\small\color{quiet}Research report / 24 November 2026}
\fancyfoot[C]{\small\thepage}
\setlength{\headheight}{14pt}
\lstset{language=Python,basicstyle=\ttfamily\scriptsize,
  breaklines=true,columns=fullflexible,keepspaces=true,
  showstringspaces=false,frame=single,rulecolor=\color{quiet!35},
  aboveskip=8pt,belowskip=8pt}
\hypersetup{pdftitle={Graffiti Transfers: Sharper Extremal Bounds and Arithmetic Rank Obstructions},
pdfauthor={Keio and Shaitetsu},pdfsubject={Research deductions with explicit dependencies and scope}}
\begin{document}
\begin{center}
{\small\color{quiet}VERIFIABLE MATHEMATICAL REASONING / RESEARCH CONTINUATION}\\[9pt]
{\LARGE\bfseries\color{ink}Graffiti Transfers}\\[6pt]
{\Large A sharper extremal bound, arithmetic counterexamples,\\[3pt]
and a Log-Rank reformulation}\\[10pt]
{\large Keio and Shaitetsu}\\[4pt]
24 November 2026
\end{center}
\vspace{3pt}\hrule\vspace{9pt}

\textbf{Abstract.}
Starting from the weighted distance theorem recorded in VMRC manuscript
[227], we lower its explicit sufficient order for the sharp
Caporossi--Hansen additive bound from $10^{27}$ to $10^{10}$.
The star is the unique minimizer throughout the new range.
The key improvement is an independently proved, exact extremal formula
for the distance of two samples on a tree with positive integer vertex
weights. We also transport finite graphs to squarefree integers and
construct an explicit family whose noncoprimality indicator matrix has
rank $29^k$, while partitioning the integers into pairwise-coprime classes
requires exactly $32^k$ classes. This transfers and amplifies a classical
Alon--Seymour counterexample; it is not a new historical conjecture
resolution. The same dictionary gives an exact arithmetic version of
the Log-Rank conjecture. Finally, applying the corpus weighted theorem
coordinate by coordinate gives a packing bound for products of graph
alphabets, including a sharp gcd-chain example.

\section{Results, dependencies, and scope}

The starting question is whether the resolved assertions in
Fajtlowicz's 1988 list can be transferred to other fields.
The original list contains statements later disproved as well as
statements proved. For example, its original item 18 is
$\chi(G)\le 1+\rank A_G$; the graph of Alon and Seymour has
$\chi=32$ and $\rank A=29$ \cite{fajtlowicz,alonseymour}.
Accordingly, a transfer must preserve whether it uses a theorem
or a counterexample.

\begin{center}
\small
\begin{tabularx}{\linewidth}{@{}>{\raggedright\arraybackslash}p{.27\linewidth}X>{\raggedright\arraybackslash}p{.24\linewidth}@{}}
\toprule
Result & Exact scope & Dependency/status\\
\midrule
Caporossi--Hansen gap &
All connected graphs of every order $n\ge10^{10}$;
equality exactly on stars &
New deduction here from the corpus weighted theorem\\
Weighted tree extremum &
Every number of edges and every admissible integer total weight &
Independent proof; sharp formula\\
Arithmetic rank obstruction &
An explicit family for every integer $k\ge1$ &
Classical graph seed, exact transfer and tensor proof\\
Arithmetic Log-Rank statement &
Equivalent to the full graph formulation &
Reformulation; unresolved\\
Product alphabet packing &
Arbitrary finite connected graph factors and positive coordinate costs &
Corollary of the corpus weighted theorem\\
\bottomrule
\end{tabularx}
\end{center}

The all-order Caporossi--Hansen conjecture is \emph{not} proved here.
Relative to the small-order verification reported in [227], the
remaining interval is $10\le n<10^{10}$. Neither the Log-Rank conjecture
nor a general prime-distribution or irrationality problem is resolved.
The exact tree estimate and the new numerical cutoff are proved in this
report; historical priority for these deductions has not been established.

\subsection{The corpus input}

Graphs are finite, simple and undirected. Metric assertions concern
connected graphs of order at least two. Write
\[
R(G)=\sum_{uv\in E(G)}\frac1{\sqrt{d_ud_v}},\qquad
M(G)=\sum_{uv\in E(G)}\frac1{\max(d_u,d_v)},\qquad
\mu(G)=\binom n2^{-1}\sum_{\{u,v\}}d_G(u,v).
\]
Here $M$ always denotes the modified Randi\'c sum, not a modal distance.
The weighted distance statement used from [226--227] is
\begin{equation}
 \sum_{u,v}p_up_vd_G(u,v)+2\sum_vp_v^2\le M(G)+1
 \quad\text{for every probability vector }p.
 \tag{W}\label{eq:W}
\end{equation}
The distances have zero diagonal; the diagonal correction is explicitly
$2\sum p_v^2$. We use (W) as the stated corpus theorem, with its proof
located in the frozen source \cite{vmrc227}. The present report supplies
the proofs of the additional deductions, not a new external
certification of every manuscript in the corpus.

Uniform mass and mass $1/2$ at two diameter endpoints give
\begin{equation}
 M(G)\ge\frac{(n-1)\mu(G)-(n-2)}n,\qquad M(G)\ge\frac{\diam(G)}2.
 \label{eq:uniform-input}
\end{equation}
We also need the following elementary identity and its consequence.
For $1\le a\le b\le\Delta$,
\[
(a^{-1/2}-b^{-1/2})^2
 =(\sqrt{b/a}-1)\bigl((ab)^{-1/2}-b^{-1}\bigr).
\]
Summing, and writing $r=\sqrt\Delta$, yields
\begin{equation}
 n-2R=\sum_{uv\in E}(d_u^{-1/2}-d_v^{-1/2})^2
 \le (r-1)(R-M),
 \quad (r+1)R\ge n+(r-1)M.
 \label{eq:bridge-input}
\end{equation}
These use no conjecture: $\sum_{uv\in E}(1/d_u+1/d_v)=n$.

% The next section is expanded into this single source at build time.

% Internal research component. Main document must define theorem/lemma environments.
% Dependencies: the corpus weighted-distance theorem, only through (CH-U).
% This improves an eventual cutoff; it does not prove the remaining finite range.

\section{A sharp weighted-tree bound and a smaller Caporossi--Hansen cutoff}
\label{sec:ch-improved-cutoff}

For a finite simple connected graph of order $n\ge2$, put
\[
 R=\sum_{uv\in E(G)}\frac1{\sqrt{d_ud_v}},\qquad
 M=\sum_{uv\in E(G)}\frac1{\max(d_u,d_v)},\qquad
 \mu=\binom n2^{-1}\sum_{\{u,v\}\subset V(G)}d_G(u,v).
\]
The only input taken from the corpus weighted-distance theorem is its uniform
specialization
\begin{equation}\label{eq:ch-uniform-input}\tag{CH-U}
 M\ge \frac{(n-1)\mu-(n-2)}n.
\end{equation}
Indeed, the corpus theorem states that
\[
 \sum_{u,v}p_up_vd_G(u,v)+2\sum_vp_v^2\le M+1
\]
for every probability vector $p$; setting $p_v=1/n$ gives
\eqref{eq:ch-uniform-input}. All further estimates used in this section are
proved below. In particular, the argument does not presume the desired
Caporossi--Hansen additive inequality.

\subsection{An exact weighted-tree extremum}

\begin{lemma}\label{lem:ch-exact-weighted-tree}
Let $T$ have $h$ edges and positive integer vertex weights $w_v$ with total
$n$. Let $X,Y$ be independent samples with probability $w_v/n$ at $v$.
Then
\begin{equation}\label{eq:ch-exact-weighted-tree}
 \mathbb E d_T(X,Y)
 \le \frac h2-\frac{h(h^2-1)}{6n^2}
       -\frac{h}{2n^2}\,\boldsymbol{1}_{\{n+h\ {
m even}\}}.
\end{equation}
For $h\ge1$, equality holds precisely when $T$ is a path, every internal
vertex has weight one, and its two endpoint weights differ by at most one.
For $h=0$ the assertion is the equality $0=0$.
\end{lemma}

\begin{proof}
An edge whose deletion leaves weights $a_e,n-a_e$ contributes
$2a_e(n-a_e)/n^2$ to the expected distance. Let
$s_e=\min(a_e,n-a_e)$.

Fix an integer $a$ with $1\le a\le\lfloor n/2\rfloor$. The edges with
$s_e\ge a$, if any, form a connected subtree $S$: every edge on the path
between two such edges separates an outer side of each of those two cuts,
and both outer sides have weight at least $a$. Suppose $S$ has $m\ge1$
edges and $\ell\ge2$ leaves. At each leaf of $S$, the component pointing
away from its unique incident edge in $S$ has weight at least $a$. These
$\ell$ components are disjoint. The other $m+1-\ell$ vertices of $S$
each have weight at least one and lie outside those components. Consequently
\[
 n\ge \ell a+(m+1-\ell)\ge m+2a-1,
 \qquad m\le n-2a+1.
\]
Order the cut sizes as $s_1\ge\cdots\ge s_h$. Taking $a=s_j$ gives
\[
 s_j\le \left\lfloor\frac{n-j+1}{2}\right\rfloor
       \quad(1\le j\le h).
\]
Since $s(n-s)$ increases on $0\le s\le n/2$, the sum of the cut products
is at most the sum of the $h$ largest cut products of the unit-weight path
on $n$ vertices. The corresponding deficits from $n/2$ are the first $h$
entries of
\[
 0,1,1,2,2,\ldots\quad(n\text{ even}),
 \qquad
 \tfrac12,\tfrac12,\tfrac32,\tfrac32,\ldots\quad(n\text{ odd}).
\]
Their squared sum equals
\[
 \frac{h(h^2-1)}{12}
   +\frac h4\,\boldsymbol{1}_{\{n+h\ {
m even}\}}.
\]
Substitute this into
\[
 \mathbb E d_T(X,Y)
   =\frac h2-\frac2{n^2}\sum_e(n/2-s_e)^2
\]
to obtain the bound.

The path in the statement attains exactly these central cut sizes, so the
bound is sharp. For completeness, equality also forces that form. Every
ordered cut size must attain its upper bound, since the integer cut-product
function is strictly increasing up to $\lfloor n/2\rfloor$. If $n$ is even,
there is a central edge cutting the total weight equally. The next required
cut sizes, one on each side, are $n/2-1$, then $n/2-2$, and so on. To obtain
a component of weight one less within a side, the intervening vertex must
have weight one and must have just one continuing branch. This forces two
chains from the central edge, of lengths differing by at most one. If $n$
is odd and $h\ge2$, the two cuts of smaller-side weight $(n-1)/2$ isolate
two disjoint components, leaving a single intervening vertex of weight one;
the same argument then forces the two chains. At the final unpaired cut
either chain may be extended. All edges are exhausted, and the remaining
endpoint weights are the two integers nearest $(n-h+1)/2$. The case $h=1$
states directly that the two endpoint weights are as equal as possible.
\end{proof}

\subsection{Two elementary graph estimates}

\begin{lemma}\label{lem:ch-bridge-and-broom}
Write $\Delta=\max_vd_v$, $r=\sqrt\Delta$, and $k=n-1-\Delta$.
Then
\begin{align}
 n-2R&\le(r-1)(R-M),\label{eq:ch-edge-bridge}\\
 \mu&\le 2-\frac2n+
 \frac{k(k+1)(3n-2k-7)}{3n(n-1)}.\label{eq:ch-broom-sharp}
\end{align}
\end{lemma}

\begin{proof}
For an edge with endpoint degrees $a\le b$,
\[
 (a^{-1/2}-b^{-1/2})^2
   =\left(\sqrt{b/a}-1\right)
      \left((ab)^{-1/2}-b^{-1}\right).
\]
The first factor lies in $[0,r-1]$. The sum of the squares on the left is
$n-2R$, because every vertex contributes its degree times the reciprocal
of its degree. Summation proves \eqref{eq:ch-edge-bridge}.

For \eqref{eq:ch-broom-sharp}, choose a vertex of degree $d=\Delta$ and a
spanning tree containing all $d$ of its incident edges. Start with that
$(d+1)$-vertex star and add the remaining tree vertices one at a time, each
as a leaf. At a step whose resulting order is $t$, let
$b_1\le\cdots\le b_{t-1}$ be the distances from the new leaf to the earlier
vertices. A path realizing distance $j$ contains all smaller positive
distances, so $b_i\le i$. A tree on $t$ vertices containing a vertex of
degree at least $d$ has diameter at most $t-d+1$: a path contains at most
two of the edges incident to a vertex on it, and counting off-path vertices
gives this bound; if the vertex is off the path, the tree condition gives
an even stronger count. Thus $b_i\le t-d+1$ as well. Hence the new
transmission is at most
\[
 \sum_{i=1}^{t-1}\min(i,t-d+1)
      =\binom t2-\binom{d-1}{2}.
\]
The initial star has unordered distance sum $d^2$. Distances in the original
graph are no greater than those of the spanning tree. Therefore
\[
 \binom n2\mu
 \le d^2+\sum_{t=d+2}^{n}
                  \left(\binom t2-\binom{d-1}{2}\right).
\]
Evaluate the sum and put $d=n-1-k$ to obtain
\eqref{eq:ch-broom-sharp}. The bound is attained by the broom obtained by
extending one arm of the initial star.
\end{proof}

\subsection{The improved eventual result}

\begin{theorem}\label{thm:ch-cutoff-ten-billion}
For every $n\ge10^{10}$, every finite simple connected graph $G$ of order
$n$ satisfying \eqref{eq:ch-uniform-input} obeys
\[
 R(G)-\mu(G)\ge\sqrt{n-1}-2+\frac2n.
\]
Equality holds precisely for the star $K_{1,n-1}$. Thus the corpus
weighted-distance theorem proves the Caporossi--Hansen additive bound, with
its unique equality graph, for all $n\ge10^{10}$.
\end{theorem}

\begin{proof}
Set
\[
 s=\sqrt{n-1},\quad r=\sqrt\Delta,\quad
 g=R-\mu,\quad \gamma=s-2+\frac2n,\quad X=R-M.
\]
It suffices to show that $g\le\gamma$ forces $G$ to be the star.
From \eqref{eq:ch-uniform-input},
\begin{equation}\label{eq:ch-X-bound}
 0\le X\le g+1+\frac{\mu-2}{n}
             \le s-1+\frac\mu n<s<\sqrt n.
\end{equation}
Here $M\le R$ follows termwise, and $\mu\le n-1<n$ suffices for the
strict bound. We also have $g\le\gamma<\sqrt n$.

\paragraph{Concentrating the graph around vertices of large degree.}
Choose
\[
 K=\lceil10n^{1/3}\rceil,\qquad t=\sqrt{K-1}<\sqrt{10}\,n^{1/6}.
\]
Call a vertex high if its degree is at least $K$, and choose a maximal
collection of $c$ high vertices at pairwise distance at least three. Their
closed radius-one balls are disjoint, giving
\begin{equation}\label{eq:ch-c-bound}
 c\le\frac n{K+1}\le\frac{n^{2/3}}{10}.
\end{equation}
Let $A$ be the union of their closed radius-three balls, and put
$b=n-|A|$. Every high vertex is within distance two of a chosen vertex,
so both ends of every edge incident to a high vertex belong to $A$.

In the sum representing $n-2R$, the edges incident to a high vertex
contribute at most $|A|$, by replacing each squared difference by
$1/d_u+1/d_v$ and summing the incident fractions at each vertex of $A$.
On every other edge both degrees are at most $K-1$, and the first factor
in the proof of \eqref{eq:ch-edge-bridge} is at most $t-1$. Thus
\begin{equation}\label{eq:ch-R-split}
 R\ge\frac b2-\frac{t-1}{2}X.
\end{equation}

Contract each connected component of $G[A]$ to a vertex of weight equal
to its order, and leave each vertex outside $A$ with weight one. Let
$c_A$ be the number of contracted components; then $c_A\le c$. Choose a
spanning tree of the connected quotient. Its number of edges is
\[
 h=b+c_A-1,\qquad 0\le h\le n-1,
 \qquad b\ge h-c+1.
\]
These identities also cover $A=\varnothing$, when $c_A=c=0$ and $h=n-1$.
A component containing $c_i$ chosen balls has internal diameter at most
$7c_i-1$: each ball has radius three, and a chain of balls can be joined
using at most one extra edge between consecutive balls. Every path in the
quotient tree consequently lifts to a walk in $G$ with at most $7c$
additional edges. This estimate also handles two endpoints in the same
contracted component.

Apply Lemma~\ref{lem:ch-exact-weighted-tree}, dropping only its nonnegative
parity correction. Two independent uniform original vertices have mean
distance $(n-1)\mu/n$. Therefore
\begin{equation}\label{eq:ch-mu-quotient}
 \mu\le\frac n{n-1}
       \left(\frac h2-\frac{h^3-h}{6n^2}+7c\right).
\end{equation}
Subtract this from \eqref{eq:ch-R-split}, using $b\ge h-c+1$ and
$h\le n-1$. Rearrangement gives
\begin{align}
 \frac{h^3}{6n(n-1)}
 &\le g+\frac{t-1}{2}X+
     \left(\frac12+\frac{7n}{n-1}\right)c+\frac1{6n}
       \label{eq:ch-cubic-before}\\
 &<\left(\frac{\sqrt{10}}2+\frac34\right)n^{2/3}
       +\frac12\sqrt n+\frac{7n^{2/3}}{10(n-1)}+\frac1{6n}
       \nonumber\\
 &<\frac73n^{2/3}+\frac12\sqrt n+1.
       \label{eq:ch-cubic-after}
\end{align}
For the first line, the terms $h/[2(n-1)]$ and $-1/2$ cancel in the
required direction, and $h/[6n(n-1)]\le1/(6n)$. For the last line,
$\sqrt{10}<19/6$ and
$7n^{2/3}/[10(n-1)]+1/(6n)<1$ at every $n\ge10^{10}$.

Consequently
\[
 \left(\frac hn\right)^3
   <14n^{-1/3}+3n^{-1/2}+\frac6n.
\]
For $n\ge10^{10}$, use $n^{-1/3}<1/2100$,
$n^{-1/2}\le1/100000$ and $n^{-1}\le10^{-10}$. The rational comparison
\[
 \frac1{150}+\frac3{100000}+\frac6{10^{10}}
        <\left(\frac{189}{1000}\right)^3
\]
shows $h/n<189/1000$. Since $h^3-h\ge0$, equations
\eqref{eq:ch-mu-quotient} and \eqref{eq:ch-c-bound} then give
\begin{equation}\label{eq:ch-mu-small}
 \frac\mu n
 <\frac n{n-1}\left(\frac{189}{2000}+\frac7{10}n^{-1/3}\right)
 <\frac{10^{10}}{10^{10}-1}
       \left(\frac{189}{2000}+\frac1{3000}\right)
 <\frac{19}{200}.
\end{equation}

\paragraph{A vertex adjacent to more than $65\%$ of the graph.}
Combining \eqref{eq:ch-edge-bridge}, $R\le\mu+\gamma$ and
\eqref{eq:ch-X-bound} gives
\[
 n-2\mu-2\gamma\le(r-1)X\le(r-1)s.
\]
Because $1-2\mu/n>0$, substitution of $\gamma$ yields
\[
 \frac r{\sqrt n}
 \ge \frac{\sqrt n}{s}\left(1-\frac{2\mu}{n}\right)
      -\frac1{\sqrt n}+\frac{4(1-1/n)}{s\sqrt n}
 >\frac{81}{100}-\frac1{100000}>\frac{809}{1000}.
\]
In particular, $(809/1000)^2>13/20$ implies
\begin{equation}\label{eq:ch-k-small}
 \frac{k}{n}=1-\frac1n-\frac{r^2}{n}<\frac7{20},
 \qquad k=n-1-\Delta.
\end{equation}

\paragraph{Eliminating every positive deficit in the maximum degree.}
Suppose $k\ge1$, and put $A_r=2+(r-1)/n$. The bridge together with
\eqref{eq:ch-uniform-input} gives
\[
 (r+1)g\ge n-2\mu-(r-1)\left(1+\frac{\mu-2}{n}\right).
\]
Using the identity
$n-r-3=(r+1)(s-2)+(s-1)(s-r)$ and $g\le\gamma$, we obtain
\[
 (s-1)(s-r)\le\frac{2(r+1)}n+A_r(\mu-2).
\]
Now apply \eqref{eq:ch-broom-sharp}, use $s-r=k/(s+r)$, and divide by
$k>0$. This forces
\begin{equation}\label{eq:ch-final-forced}
 \frac{s-1}{s+r}
 \le\frac{2(r-1)(n-1)}{n^2k}
   +A_r\frac{(k+1)(3n-2k-7)}{3n(n-1)}.
\end{equation}

Let $x=k/n$. Throughout $0<x<7/20$ by \eqref{eq:ch-k-small}. The right
side of \eqref{eq:ch-final-forced} has the elementary upper bounds
\[
 \frac{2(r-1)(n-1)}{n^2k}\le\frac2{\sqrt n}\le\frac1{50000},
 \qquad A_r\le2+\frac1{100000},
\]
and
\begin{equation}\label{eq:ch-final-fraction}
 \frac{(k+1)(3n-2k-7)}{3n(n-1)}
 \le x-\frac23x^2+\frac2n.
\end{equation}
To check the last bound, write $\varepsilon=1/n$ and
$f(x)=x-2x^2/3$. The left side is
$(x+\varepsilon)(1-2x/3-7\varepsilon/3)/(1-\varepsilon)$,
at most $(f(x)+\varepsilon)/(1-\varepsilon)$.
Here $f(x)\le f(7/20)=161/600$ and $\varepsilon\le10^{-10}$, so this
is at most $f(x)+2\varepsilon$.

If $x\le3/10$, the right side of \eqref{eq:ch-final-forced} is at most
\[
 \frac1{50000}+
 \left(2+\frac1{100000}\right)
 \left(\frac6{25}+\frac2{10^{10}}\right)
 <\frac{481}{1000}.
\]
But its left side is at least $(s-1)/(2s)>499/1000$, a contradiction.

If $3/10<x<7/20$, then $r<\sqrt{7n/10}$ and $s-1\ge\sqrt n-2$.
Since $\sqrt{7/10}<21/25$, its left side is greater than
\[
 \frac{1-2/100000}{1+21/25}>\frac{27}{50}.
\]
Its right side is at most
\[
 \frac1{50000}+
 \left(2+\frac1{100000}\right)
 \left(\frac{161}{600}+\frac2{10^{10}}\right)
 <\frac{537}{1000}<\frac{27}{50},
\]
another contradiction. Hence $k=0$.

\paragraph{Equality.}
The graph therefore has a dominating vertex. Its incident edges contribute
exactly one to $M$, so $M\ge1$. With $r=s$, the bridge gives
\[
 R\ge\frac{n+(s-1)M}{s+1}
       \ge\frac{n+s-1}{s+1}=s.
\]
All pair distances are at most two, so
$\mu=2-|E(G)|/\binom n2\le2-2/n$. If any edge beyond the spanning star
is present, this last inequality is strict. Thus $g\ge\gamma$, with
equality precisely on the star. Since the proof started from $g\le\gamma$,
it also excludes every strict counterexample.
\end{proof}

\paragraph{Exact scope of the improvement.}
This argument replaces the earlier sufficient cutoff $10^{27}$ by
$10^{10}$, while retaining the same corpus weighted-distance input.
Together with the previously recorded verification for $2\le n\le9$, the
remaining interval for a universal proof is exactly
\[
                    10\le n<10^{10}.
\]
No claim of an all-order resolution is made by this section.

\section{Exact transport from graphs to arithmetic}

For a finite set $S=\{a_1,\ldots,a_n\}$ of distinct positive integers,
define its \emph{noncoprimality indicator matrix} by
\[
C(S)_{ij}=
\begin{cases}
1,&i\ne j\text{ and }\gcd(a_i,a_j)>1,\\
0,&\text{otherwise}.
\end{cases}
\]
It is important that this is a zero-diagonal $0$--$1$ matrix,
not the usual gcd-value matrix $(\gcd(a_i,a_j))$.
Write $\rho(S)=\rank_{\R}C(S)$, and let $\cchi(S)$ be the smallest
number of classes in a partition of $S$ such that every two distinct
integers in the same class are relatively prime.
These are respectively adjacency rank and chromatic number of the
noncoprimality graph.

\begin{theorem}[Squarefree representation]\label{thm:representation}
Every finite simple graph $G=(V,E)$ has distinct squarefree integer
labels $(a_v)_{v\in V}$ such that
\[
uv\in E\quad\Longleftrightarrow\quad \gcd(a_u,a_v)>1
\qquad(u\ne v).
\]
Every nontrivial gcd can be required to be a single prime, and
no prime need divide more than two labels. The labels can also all
be required to have exactly $\Delta(G)+1$ distinct prime factors.
\end{theorem}
\begin{proof}
Choose a distinct prime $p_e$ for every edge. At vertex $v$ choose
$\Delta+1-d_v$ additional private primes, all disjoint from the edge
primes and from each other's private primes. Let $a_v$ be the product
of its incident edge primes and its private primes. A prime common to
two labels can only be the prime on their common edge. Each label has
at least one private prime, so the labels are distinct. All the asserted
properties follow.
\end{proof}

This is a prime-product form of the classical representation of graphs
by intersections of finite sets \cite{egp}. No originality is claimed
for the underlying representation. Its value here is the exact
preservation of all relevant parameters.

\begin{center}\small
\begin{tabularx}{\linewidth}{@{}lX@{}}
\toprule
Graph parameter & Arithmetic meaning\\
\midrule
Independence number & Largest pairwise-coprime subfamily\\
Chromatic number & Fewest pairwise-coprime classes in a partition\\
Adjacency rank & Real rank of $C(S)$\\
Vertex degree & Number of other integers sharing a prime factor\\
Graph distance & Shortest chain of nontrivial gcds\\
\bottomrule
\end{tabularx}
\end{center}

\subsection{A direct arithmetic consequence of items 13 and 14}

Suppose the noncoprimality graph of $S$ is connected. Put
$d(a)=|\{b\in S\setminus\{a\}:\gcd(a,b)>1\}|$ and let
$\ell(a,b)$ be its gcd-chain distance. The corpus mean-distance bound
translates exactly into
\[
\binom{|S|}{2}^{-1}\sum_{\{a,b\}\subset S}\ell(a,b)
\le
\sum_{\substack{\{a,b\}\subset S\\\gcd(a,b)>1}}
\frac1{\sqrt{d(a)d(b)}}.
\]
Chung's inequality from original item 14 also guarantees a
pairwise-coprime subfamily of size at least the ceiling of the
left-hand side \cite{fajtlowicz}.
These are transfer corollaries; they do not assert new facts about
the distribution of prime numbers.

\subsection{An elementary strengthening of the rank input}

\begin{proposition}\label{prop:rank-density}
For every nonempty finite simple graph,
\[
\rank A_G\ \ge\ \omega(G)\
\ge\ \sum_{v\in V(G)}\frac1{n-d_v}.
\]
Consequently, if $C(S)\ne0$, then
$\sum_{a\in S}1/(|S|-d(a))\le\rho(S)$.
\end{proposition}
\begin{proof}
Order the vertices of the complement uniformly at random and select
each vertex that precedes all its complement-neighbours.
The selected vertices are independent in the complement, hence a
clique in $G$. Vertex $v$ is selected with probability
$1/(n-d_v)$, proving the second inequality by expectation.
A maximum clique has size at least two, and its adjacency matrix is
$J-I$, with eigenvalues $\omega-1,-1,\ldots,-1$. It is nonsingular.
As a principal submatrix of $A_G$, its rank cannot exceed $\rank A_G$.
\end{proof}

The original average-temperature expression is
$T(G)=n^{-1}\sum_vd_v/(n-d_v)=\sum_v1/(n-d_v)-1$.
Thus Proposition~\ref{prop:rank-density} is already stronger than
original item 1. It is the classical random-order independent-set bound
applied to the complement, not a new rank theorem.

\section{An infinite arithmetic rank obstruction}

\begin{theorem}\label{thm:arithmetic-family}
For every integer $k\ge1$ there is an explicit finite set $S_k$ of
distinct squarefree positive integers such that
\[
|S_k|=64^k,\qquad \rho(S_k)=29^k,\qquad \cchi(S_k)=32^k.
\]
Each integer has exactly $64^k-8^k$ distinct prime factors.
Every prime used divides exactly two integers, and the gcd of any
two distinct members is either $1$ or a single prime.
\end{theorem}

\subsection{The classical 64-vertex seed}

Let $V=\F_2^6$. Join $x,y$ when $2\le\wt(x+y)\le5$.
This is the Alon--Seymour graph \cite{alonseymour}.
Its complement has difference set
$\{e_1,\ldots,e_6,\one\}$. The sum of two distinct members of this
set has weight $2$ or $5$, so the complement is triangle-free.
It has edges; hence $\alpha(G)=2$.
Color $x$ by its last five coordinates. Each color class has two
nonadjacent vertices. Therefore $\chi(G)=32$.

Put $B=J_{64}-A_G$. Its Walsh character indexed by $s\in\F_2^6$
has eigenvalue
$\lambda_j=1+6-2j+(-1)^j$, where $j=\wt(s)$.
Indeed, the eigenvalue is the character sum over
$\{0,e_1,\ldots,e_6,\one\}$.

\begin{center}
\begin{tabular}{@{}crrrrrrr@{}}
\toprule
$j$ &0&1&2&3&4&5&6\\
\midrule
$\lambda_j$ &8&4&4&0&0&$-4$&$-4$\\
Multiplicity &1&6&15&20&15&6&1\\
\bottomrule
\end{tabular}
\end{center}
Thus $B$ has rank $29$. On the constant character $A_G$ has
eigenvalue $56$; on other characters it has eigenvalue
$-\lambda_j$. In particular,
\[
\operatorname{spec}(A_G)=
\{56^{(1)},\,(-4)^{(21)},\,0^{(35)},\,4^{(7)}\},\qquad \rank A_G=29.
\]

\subsection{Tensor amplification with an exact rank count}

On $V^k$ join two tuples if at least one coordinate pair is an edge
of $G$. The resulting simple graph has adjacency matrix
\[
A_k=J_{64^k}-B^{\otimes k}.
\]
The diagonal is zero because all diagonal entries of $B$ are one.
The tensor Walsh basis diagonalizes $A_k$. The constant eigenvalue is
$64^k-8^k\ne0$. All other eigenvalues are the negatives of products of
the $k$ corresponding $B$-eigenvalues. There are $29^k$ choices with
all factors nonzero, including the constant one. Therefore
\[
\rank A_k=29^k.
\]
An independent set projects in each coordinate to an independent set
of $G$, of size at most two. It therefore has size at most $2^k$,
and the product of $k$ independent pairs attains this bound.
The product coloring uses $32^k$ colors, while
$|V^k|/\alpha(G_k)=32^k$ proves the reverse inequality. Hence
\[
\alpha(G_k)=2^k,\qquad\chi(G_k)=32^k.
\]

\subsection{Integer labels and the false universal assertion}

Assign the first $|E(G_k)|$ primes to the edges in lexicographic order.
For every vertex $v$, define
\[
a_v=\prod_{e\ni v}p_e.
\]
The graph is regular of degree $d_k=64^k-8^k\ge56$. Distinct vertices
share at most one incident edge, so their prime sets and their labels
are distinct. Every prime occurs twice. Moreover $C(S_k)=A_k$ for
$S_k=\{a_v:v\in V^k\}$. This proves
Theorem~\ref{thm:arithmetic-family}.

\begin{center}
\begin{tabular}{@{}rrrr@{}}
\toprule
$k$ & Number of integers & $\rho(S_k)$ & $\cchi(S_k)$\\
\midrule
1&64&29&32\\
2&4,096&841&1,024\\
3&262,144&24,389&32,768\\
\bottomrule
\end{tabular}
\end{center}

\begin{corollary}
There is no absolute constant $C$ such that
$\cchi(S)\le C\rho(S)$ for every finite integer set $S$, even under
the squarefree and prime-sharing restrictions above.
More generally $\cchi(S)\le C\rho(S)^\beta$ fails for every
$\beta<\log32/\log29=1.0292341623\ldots$.
\end{corollary}
\begin{proof}
The ratio in the first assertion is $(32/29)^k\to\infty$.
For the second it is $(32/29^\beta)^k\to\infty$.
\end{proof}

This closes a precise proposed \emph{linear arithmetic principle} by a
full counterexample family. It does not constitute a new disproof of
the historical rank-coloring conjecture: the seed is classical, and
stronger graph rank--coloring separations are already known
\cite{razborov,cioabatait}.

\section{A major equivalent target: arithmetic Log-Rank}

The graph formulation of the Log-Rank conjecture asks whether there
are absolute constants $C,c>0$ such that
\begin{equation}
\log\chi(G)\le C\bigl(\log(\rank A_G+2)\bigr)^c
\quad\text{for every finite simple graph }G.
\label{eq:graph-logrank}
\end{equation}
The harmless $+2$ covers small or zero rank. Its equivalence to the
communication-complexity formulation is established in the literature;
recent sources continue to treat the conjecture as open
\cite{rankramsey,ccc2026}.

\begin{proposition}[Arithmetic equivalence]\label{prop:logrank}
Equation~\eqref{eq:graph-logrank} is equivalent to the existence of
absolute $C,c>0$ such that
\[
\cchi(S)\le
\exp\!\left(C\bigl(\log(\rho(S)+2)\bigr)^c\right)
\quad\text{for every finite integer set }S.
\]
It is enough to establish the arithmetic assertion on squarefree sets
with uniform numbers of prime factors, at most two uses of each prime,
and pairwise gcds equal to $1$ or a prime.
\end{proposition}
\begin{proof}
Apply the graph assertion to the noncoprimality graph of $S$.
Conversely, represent an arbitrary graph by
Theorem~\ref{thm:representation}; both its rank and its chromatic
number are preserved. The restricted version suffices because that
representation has all the stated restrictions.
\end{proof}

The family of Theorem~\ref{thm:arithmetic-family} does not disprove
Log-Rank: it exhibits a fixed polynomial separation. On this family,
the rank-density lower bound is only
\[
\sum_v\frac1{64^k-d_k}=8^k,
\]
whereas the rank is $29^k$ and the chromatic number is $32^k$.
The bound controls clique size or density, not a sufficiently small
coloring. A new upper bound on partition complexity is still required.

\section{Coordinatewise transfer to packing and coding}

\subsection{One graph: separated sets and collision probability}

Let $C\subseteq V(G)$ have size $K$ and pairwise distances at least
$\delta$. Put uniform mass on $C$ in (W). Then
\[
\delta(1-1/K)+2/K\le M(G)+1.
\]
If $\delta>M(G)+1$, this gives
\begin{equation}
K\le\left\lfloor\frac{\delta-2}{\delta-M(G)-1}\right\rfloor.
\label{eq:single-packing}
\end{equation}
For any probability vector supported on such a separated set, the
collision probability $Q=\sum p_v^2$ satisfies
\[
Q\ge\frac{\delta-M(G)-1}{\delta-2}.
\]
There is no sign issue in these rearrangements: $M(G)\ge1$, so the
assumption implies $\delta>2$.
For a path of length $D>2$ and its two endpoints, $M=D/2$ and
\eqref{eq:single-packing} is attained with equality.

\subsection{Products without an exponential graph budget}

\begin{theorem}[Coordinatewise packing bound]\label{thm:product-packing}
Let $G_i$ be finite connected graphs, let
$S_i\subseteq V(G_i)$ have size $s_i\ge1$, and let $c_i>0$.
For a code $C\subseteq\prod_{i=1}^t S_i$ of size $K$ use the metric
$d(x,y)=\sum_i c_i d_{G_i}(x_i,y_i)$.
If its minimum distance is $\delta$, put
\[
B_*=\sum_{i=1}^t c_i\left(M(G_i)+1-\frac2{s_i}\right).
\]
Then
\[
\delta(1-1/K)\le B_*.
\]
In particular, if $\delta>B_*$, then
$K\le\lfloor\delta/(\delta-B_*)\rfloor$.
\end{theorem}
\begin{proof}
Choose $X,Y$ independently and uniformly from the code.
Then $\E d(X,Y)\ge\delta(1-1/K)$.
Let $p_i$ be the $i$th marginal distribution. Apply (W) to each factor:
\[
\E d_{G_i}(X_i,Y_i)
\le M(G_i)+1-2\sum_{a\in S_i}p_i(a)^2
\le M(G_i)+1-2/s_i.
\]
The last step is Cauchy--Schwarz. Multiply by $c_i$ and sum.
\end{proof}

Keeping the actual marginal counts $n_{i,a}$ gives the stronger exact
necessary condition
\[
\delta K(K-1)\le\sum_i c_i
\left[(M(G_i)+1)K^2-2\sum_a n_{i,a}^2\right].
\]
For rational $c_i$, all denominators can be cleared, giving an
integer-arithmetic certificate.

\subsection{A sharp arithmetic example}

For each coordinate use the integer graph $\{2,3,6\}$, with the path
$2-6-3$, and allow symbols only from $\{2,3\}$.
Here $M=1$, $s_i=2$, and the two symbols have gcd-chain distance two.
For three coordinates $B_*=3$. The four tuples
\[
(2,2,2),\quad(2,3,3),\quad(3,2,3),\quad(3,3,2)
\]
have minimum distance four. Theorem~\ref{thm:product-packing} gives
$K\le4/(4-3)=4$, so the optimum is exactly four.

Applying the single-graph bound directly to the full Cartesian product
$P_3^{\square3}$ would give
\[
M(P_3^{\square3})=\frac{24}4+\frac{24}5+\frac66=\frac{59}5.
\]
Its threshold $M+1=64/5$ exceeds the distance four under consideration.
The coordinatewise application preserves information that is lost in
the full-graph budget.

For star factors with $q$ allowed leaves this reproduces the classical
$q$-ary Plotkin bound. The sharp four-word example is the standard
binary even-parity code written in arithmetic symbols
\cite{plotkin}. These special cases are established mathematics,
not newly solved coding conjectures.

\section{What a valid closure of another problem requires}

Suppose a corpus theorem is $\forall G\,[H(G)\Rightarrow Q(G)]$.
To prove a new statement $P$, it suffices to construct, from a
hypothetical counterexample to $P$, a graph satisfying $H$ and
violating $Q$:
\[
\neg P\quad\Longrightarrow\quad
\exists G\,\bigl(H(G)\ \text{and}\ \neg Q(G)\bigr).
\]
For instance, a number-theoretic counterexample would be excluded if
it forced a connected graph with $R(G)-\mu(G)<0$.
Full two-way equivalence is unnecessary for such a proof.

A map that always produces a valid graph, followed merely by testing
a universally true graph inequality, supplies a constant affirmative
answer. The target-specific work lies in proving the incompatible
bounds or preserving an equality condition. The squarefree
representation by itself supplies no quantitative information about
the numerical height of the integers or the distribution of primes.

There can be no effective procedure settling \emph{every}
integer-polynomial solvability question by reducing it to finitely
decidable graph tests: that would contradict the undecidability of
Hilbert's tenth problem \cite{matiyasevich}. This is a limitation on a
universal decision algorithm, not on research into particular problems.

\subsection{The two remaining frontiers}

For the Caporossi--Hansen additive conjecture, the new theorem removes
all orders $n\ge10^{10}$ from the outstanding range, using the same
weighted corpus input as [227]. A proof for every order in the remaining
interval has not been obtained here.

For the arithmetic Log-Rank formulation, exact encoding is complete.
What remains is a quasi-polynomial partition upper bound. The six
affirmative graph inequalities from the original list do not themselves
supply that bound. The explicit rank family is a useful test case for
any proposed strengthening.

\section{Verification and provenance}

The corpus snapshot used was
\texttt{52c1e4ae1f82bd3fdc8b809b9a94e2e4c6ef4cb3}.
The relevant manuscripts were [226], [227], [243] and [244].
The comparison with the previous $10^{27}$ cutoff refers specifically
to Section 6 of [227], not to an exhaustive literature priority claim.

The supplemental standard-library Python script constructs the seed
graph and all 64 integer labels, retaining every prime factor and
decimal integer. It verifies all 64 Walsh eigenvectors against both
$A$ and $J-A$, 4,096 Walsh inner products, all 41,664 vertex triples for
the complement triangle check, the explicit 32-coloring, and all
2,016 pairwise gcds. The labels use 1,792 primes, each exactly twice;
every label has 56 prime factors.

An independent finite check of the weighted tree formula covered all
1,441 labelled trees on two through six vertices and all positive
integer weight assignments of total at most nine: 126,918 cases.
It also checked the 30 indicated path extremizers.
All tests passed. These finite checks are supporting diagnostics.
The all-order and all-$k$ conclusions rest on the proofs above,
not on extrapolation from these checks.

\begin{thebibliography}{99}
\bibitem{fajtlowicz}
S. Fajtlowicz, On conjectures of Graffiti,
\emph{Discrete Mathematics} \textbf{72} (1988), 113--118.
\href{https://doi.org/10.1016/0012-365X(88)90199-9}{doi:10.1016/0012-365X(88)90199-9}.
The original open list is on p.~115.

\bibitem{vmrc227}
Verifiable Mathematical Reasoning Corpus,
\emph{A proof audit of Fajtlowicz's average-distance conjecture},
manuscript [227], 21 November 2026.
\href{https://github.com/QEEDDD/Verifiable-Mathematical-Reasoning-Corpus/blob/52c1e4ae1f82bd3fdc8b809b9a94e2e4c6ef4cb3/\%5B227\%5D\%20Fajtlowicz_Proof_Audit.tex}{Frozen TeX source}.
The weighted statement is (W); the previous explicit cutoff is in
the section on the Caporossi--Hansen gap.

\bibitem{favaron}
O. Favaron, M. Mah\'eo and J.-F. Sacl\'e,
The Randi\'c index and other Graffiti parameters of graphs,
\emph{MATCH} \textbf{47} (2003), 7--23, Conjecture 4.8.
\url{https://match.pmf.kg.ac.rs/electronic_versions/Match47/match47_7-23.pdf}.

\bibitem{egp}
P. Erd\H{o}s, A. W. Goodman and L. P\'osa,
The representation of a graph by set intersections,
\emph{Canadian Journal of Mathematics} \textbf{18} (1966), 106--112.
\href{https://doi.org/10.4153/CJM-1966-014-3}{doi:10.4153/CJM-1966-014-3}.

\bibitem{alonseymour}
N. Alon and P. D. Seymour,
A counterexample to the rank-coloring conjecture,
\emph{Journal of Graph Theory} \textbf{13} (1989), 523--525.
\href{https://doi.org/10.1002/jgt.3190130413}{doi:10.1002/jgt.3190130413}.

\bibitem{razborov}
A. A. Razborov,
The gap between the chromatic number of a graph and the rank of its
adjacency matrix is superlinear,
\emph{Discrete Mathematics} \textbf{108} (1992), 393--396.
\href{https://doi.org/10.1016/0012-365X(92)90691-8}{doi:10.1016/0012-365X(92)90691-8}.

\bibitem{cioabatait}
S. M. Cioab\u{a} and M. Tait,
More counterexamples to the Alon--Saks--Seymour and rank-coloring conjectures,
\emph{Electronic Journal of Combinatorics} \textbf{18} (2011), P26.
\href{https://doi.org/10.37236/513}{doi:10.37236/513}.

\bibitem{rankramsey}
G. Beniamini, N. Linial and A. Shraibman,
The Rank-Ramsey problem and the Log-Rank conjecture,
\emph{Combinatorics, Probability and Computing} \textbf{35} (2026), 568--597.
\href{https://doi.org/10.1017/S0963548326100455}{doi:10.1017/S0963548326100455}.

\bibitem{ccc2026}
L. Hambardzumyan, S. Lovett and M. Shirley,
The Log-Rank Conjecture: New Equivalent Formulations,
\emph{CCC 2026}, LIPIcs \textbf{383}, Article 7.
\href{https://doi.org/10.4230/LIPIcs.CCC.2026.7}{doi:10.4230/LIPIcs.CCC.2026.7}.

\bibitem{plotkin}
M. Plotkin, Binary codes with specified minimum distance,
\emph{IRE Transactions on Information Theory} \textbf{6} (1960), 445--450.
\href{https://doi.org/10.1109/TIT.1960.1057584}{doi:10.1109/TIT.1960.1057584}.

\bibitem{matiyasevich}
Yu. V. Matiyasevich,
Diophantine representation of enumerable predicates,
\emph{Mathematics of the USSR-Izvestiya} \textbf{5} (1971), 1--28.
\href{https://doi.org/10.1070/IM1971v005n01ABEH001004}{doi:10.1070/IM1971v005n01ABEH001004}.
\end{thebibliography}
\end{document}
